\documentclass[../main.tex]{subfiles}
\begin{document}
\chapter{Introducción}
% Borrador solo en este capítulo (activar DESPUÉS del chapter para asegurar página correcta)
\draftwatermarkon

\todo{Completar la introducción}

% Use a refsection only when compiling the chapter standalone
\ifSubfilesClassLoaded{%
	\begin{refsection}%
}{}%

Texto de ejemplo para el capítulo 1, incluyendo una cita de ejemplo \cite{latexcompanion}.

\lipsum[1-2]

El entorno \texttt{destacado} resalta un párrafo con una barra vertical al margen
izquierdo, útil para una idea o un resultado que conviene señalar sin recargar la página:

\begin{destacado}
\lipsum[3][1-4]
\end{destacado}

\lipsum[4-14]

\section{Motivación}
\todo[inline]{Redactar la motivación del trabajo: contexto del problema, brecha que se busca cubrir y aporte de la tesis.}
\missingfigure{Agregar figura de motivación}

% Al compilar el capítulo por separado, imprimimos la bibliografía local
\ifSubfilesClassLoaded{%
	\printbibliography[title={Bibliografía del Capítulo}]%
	\end{refsection}%
}{}%
\end{document}
